\chapter{بندهای متمم و رویه‌های اجرایی}
\label{app:clauses}

این بخش شامل جزئیات فنی و رویه‌هایی است که پایداری عملیاتی مجلس مهستان را تضمین می‌کند.

\section{بندهای متمم (۲۲ تا ۳۱)}

\begin{modernbox}{بند ۲۲: نظارت بین‌المللی}
تمامی مراحل انتخابات مجلس مهستان و رفراندوم‌های بعدی باید تحت نظارت ناظران بین‌المللی (سازمان ملل، اتحادیه اروپا و نهادهای تخصصی مانند IDEA) انجام شود.
\end{modernbox}

\begin{modernbox}{بند ۲۵: شفافیت مالی و تضاد منافع}
نمایندگان مجلس مهستان و اعضای کابینه موظف به اعلام عمومی تمامی دارایی‌های خود پیش از شروع به کار هستند. پذیرش هرگونه هدیه یا رانت در دوران مسئولیت منجر به عزل فوری می‌شود.
\end{modernbox}

\begin{modernbox}{بند ۲۸: حفاظت از اسناد ملی}
در روزهای اول گذار، کمیته‌ای ویژه مأمور حفاظت از تمامی آرشیوهای دولتی، امنیتی و مالی برای جلوگیری از امحای شواهد توسط عوامل نظام سابق می‌شود.
\end{modernbox}

\section{مکانیزم عزل نماینده (Recall)}

برای اولین بار در تاریخ ایران، مکانیزم عزل نماینده پیش‌بینی شده است:

\begin{center}
\begin{tikzpicture}[node distance=2.5cm, auto]
    \node (1) [circle, draw, fill=blue!10, text width=2cm, text centered] {۱. جمع‌آوری امضا (۲۰٪ حوزه)};
    \node (2) [circle, draw, fill=blue!20, text width=2cm, text centered, right of=1, node distance=4cm] {۲. تأیید توسط آمبودزمان};
    \node (3) [circle, draw, fill=blue!30, text width=2cm, text centered, right of=2, node distance=4cm] {۳. رأی‌گیری مجدد در حوزه};
    \node (4) [circle, draw, fill=green!20, text width=2cm, text centered, below of=2, node distance=2.5cm] {عزل قطعی};

    \path [draw, -Latex] (1) -- (2);
    \path [draw, -Latex] (2) -- (3);
    \path [draw, -Latex] (3) -- (4);
\end{tikzpicture}
\end{center}

\section{جدول رویه‌های اضطراری}

\begin{center}
\begin{longtable}{R{4cm} R{5cm} R{6cm}}
\caption{رویه‌های حاکم در شرایط بحران} \\
\hline
\textbf{موقعیت} & \textbf{مرجع تصمیم‌گیر} & \textbf{محدودیت زمانی} \\
\hline
\endfirsthead
\hline
\textbf{موقعیت} & \textbf{مرجع تصمیم‌گیر} & \textbf{محدودیت زمانی} \\
\hline
\endhead
\hline
\foot
ناامنی گسترده & ستاد مشترک با تأیید مجلس & حداکثر ۷۲ ساعت بدون تأیید مجدد \\
بحران پولی & بانک مرکزی و وزیر اقتصاد & گزارش هفتگی به کمیسیون اقتصاد \\
اختلاف بین قوا & هیأت میانجی‌گری ملی & ۷ روز برای ارائه‌ی راه حل نهایی \\
\hline
\end{longtable}
\end{center}
